\documentclass{article}

\usepackage[margin=1in]{geometry}
\usepackage{lmodern}
\usepackage{amsmath,amssymb,amsthm}
\usepackage{booktabs}
\usepackage{hyperref}

\newtheorem{theorem}{Theorem}
\newtheorem{lemma}[theorem]{Lemma}
\newtheorem{proposition}[theorem]{Proposition}
\theoremstyle{definition}
\newtheorem{definition}[theorem]{Definition}
\theoremstyle{remark}
\newtheorem{remark}[theorem]{Remark}

\newcommand{\Ftwo}{\mathbb{F}_2}
\newcommand{\ttw}[1]{\texttt{#1}}

\title{A Universal Six-Symbol 2-Tag System with Cycle Detection}
\author{Lancelot de Ferri\`ere\thanks{The six-symbol construction, proofs, and Lean~4 formalization were developed with GPT~6.2~Astra.}}
\date{}

\begin{document}
\maketitle

\begin{abstract}
We present a fixed six-symbol 2-tag system over the alphabet $\{\ttw{I},\ttw{D},\ttw{z},\ttw{e},\ttw{a},\ttw{b}\}$ with total appendant length $18$ that is computationally universal under finite starting queues, with halting observed as exact queue repetition.
For any deterministic counter machine with a fresh entry state, a terminal halt state, and an enabled instruction at every nonhalting configuration, a computable encoder maps the machine and its input to a finite queue that eventually repeats identically if and only if the machine halts.
The source machine need not be reversible: a reversible front end is composed with a parity-controlled tag compiler and a forward--backward bouncing construction.
The entire reduction---front end, tag compiler, and equivalence proof---is formally verified in Lean~4, leaving only Minsky's universality theorem for counter machines as an external premise.
To our knowledge, this improves on the 19-symbol universal 2-tag system UT19 to give the smallest known universal 2-tag system, and under rule-count and transition-table-size measures yields one of the smallest known universal machines operating on finite initial configurations.
We also establish a restricted lower bound showing that four clock symbols are necessary for this one-sweep parity-clock architecture.
\end{abstract}

\section{Introduction}

A \emph{2-tag system}~\cite{Post1943} consists of a finite alphabet $\Sigma$ and a production map $P:\Sigma \to \Sigma^*$.
At each step, if the queue $q \in \Sigma^*$ has length at least $2$, the system reads the first letter $x$ of $q$, deletes the first two letters, and appends $P(x)$ to the right end.
A queue of length less than $2$ has no successor.

We fix the six-symbol alphabet $\Sigma = \{\ttw{I},\ttw{D},\ttw{z},\ttw{e},\ttw{a},\ttw{b}\}$ and the six productions
\begin{equation}\label{eq:productions}
\begin{aligned}
\ttw{I} &\to \ttw{ezII}, & \ttw{D} &\to \ttw{DDez}, & \ttw{z} &\to \ttw{e}, \\
\ttw{a} &\to \ttw{b},    & \ttw{b} &\to \ttw{aaaeaeae}, & \ttw{e} &\to \varepsilon.
\end{aligned}
\end{equation}

\begin{theorem}[Main result]\label{thm:main}
There is a computable encoding of every deterministic counter machine $M$ (having a fresh entry state, a terminal halt state, and an enabled transition at every nonhalting configuration) and input register file $x$ into a finite starting queue $q(M,x) \in \Sigma^*$ such that $M$ reaches its halt state on input $x$ if and only if the run of~\eqref{eq:productions} from $q(M,x)$ eventually repeats a queue word identically.
\end{theorem}

Here \emph{exact repetition} means there exist distinct step counts $t_1 < t_2$ at which the tag system is defined and produces the same word $q_{t_1} = q_{t_2}$; halting on a short queue does not satisfy this condition.
The source machine $M$ need not be reversible, and both program and data reside in the finite queue $q(M,x)$ while the table~\eqref{eq:productions} stays fixed.

\paragraph{Context and comparison.}
Tag systems were introduced by Post~\cite{Post1943}, proved recursively unsolvable by Minsky~\cite{Minsky1961} (deletion number $6$), and shown universal at the minimal nontrivial deletion number $2$ by Cocke and Minsky~\cite{CockeMinsky1964} (see also Wang~\cite{Wang1963}).
At the lower end of the alphabet-size spectrum, binary ($2$-symbol) $2$-tag systems are decidable---a claim of Post later proved by De~Mol~\cite{DeMol2010}---whereas $3$ symbols already suffice to simulate the Collatz $3n+1$ problem via De~Mol's system $\ttw{a} \to \ttw{bc}$, $\ttw{b} \to \ttw{a}$, $\ttw{c} \to \ttw{aaa}$~\cite{DeMol2008}.
Combined with Theorem~\ref{thm:main}, this narrows the open frontier for universal $2$-tag systems to alphabets of size $3$, $4$, and $5$.
Neary and Woods later showed that 2-tag systems can simulate Turing machines in polynomial time~\cite{NearyWoods2006} and used tag-like systems to construct small universal Turing machines~\cite{NearyWoods2009}, following Rogozhin's benchmarks~\cite{Rogozhin1996}.

Most classical universality results for 2-tag systems allow the alphabet size to grow with the simulated machine, rather than fixing a single \emph{universal} 2-tag system whose finite queue encodes both program and input.
For a fixed universal 2-tag system, the smallest previously known construction to our knowledge is \emph{UT19} (19 symbols), designed by ais523~\cite{ais523UT19}.
The table~\eqref{eq:productions} reduces the alphabet from $19$ to $6$ symbols.

Under simple rule-count and description-size metrics,~\eqref{eq:productions} is also among the smallest known universal machines of any kind: it has $6$ productions and a total right-hand-side length of $18$ symbols (counting $\varepsilon$ as $0$), compared with $22$--$25$ transitions in the smallest standard universal Turing machines such as Rogozhin's $(4,6)$ and Neary--Woods's $(6,4)$ and $(5,5)$ machines~\cite{Rogozhin1996,NearyWoods2009}.
Moreover, unlike weakly universal cellular automata such as Rule~110~\cite{Cook2004} or weakly universal small Turing machines, every computation starts from a \emph{finite} configuration with no infinite periodic background.
Two trade-offs should be noted explicitly: halting is observed by exact queue repetition rather than queue exhaustion, and counters are stored in unary-of-exponentials (giving exponential slowdown relative to the counter machine).

\paragraph{Overview of the construction.}
Because a 2-tag system deletes two letters per step, one full pass over the queue (a \emph{sweep}) reads only alternate letters---either even positions (lane $0$) or odd positions (lane $1$).
The construction exploits this alignment bit as a 1-bit control wire:
\begin{enumerate}
\item \emph{Program clocks} built from $\{\ttw{I},\ttw{D},\ttw{z},\ttw{e}\}$ sit between counters and XOR a fixed, periodic bitmask onto the passing alignment without retaining or smearing parity changes from upstream counters.
\item \emph{Exponential counters} built from $\{\ttw{a},\ttw{b}\}$ represent an exponent $E$ as $2^E$ copies of $W = \ttw{aaaeaeae}$, incrementing or decrementing $E$ in two sweeps according to the incoming lane.
\item \emph{Zero testing} uses parity underflow: decrementing at $E=0$ leaves a single odd-length letter $\ttw{b}$, which regenerates $W$ while toggling the alignment for the next command.
\item \emph{Cycle detection} uses a reversible front end and a forward--backward bouncing wrapper so that halting runs stay bounded (forcing exact repetition) while nonhalting runs grow without bound.
\end{enumerate}

\section{Sweeps and the Program Clock}

\subsection{Sweeps and alignment}

Write a state as a \emph{nominal word} $w = w^{(1)} w^{(2)} \cdots w^{(k)}$ (a concatenation of blocks) together with a \emph{cut bit} $p \in \{0,1\}$.
The corresponding \emph{physical queue} is $w$ with its first $p$ letters removed.
In one sweep, each block $u$ entered on lane $b \in \{0,1\}$ outputs the concatenation of $P(u_i)$ over zero-based indices $i \in \{b, b+2, \ldots\}$ and passes the outgoing lane $b \oplus (|u| \bmod 2)$ to the next block.
The outgoing lane after the final block becomes the next sweep's cut bit $p'$.

\begin{lemma}[Sweep semantics]\label{lem:sweep}
Suppose the leading block $w^{(1)}$ produces a nonempty word on either lane.
Then one nominal sweep corresponds to exactly $\lceil |q|/2 \rceil$ literal 2-tag steps from the physical queue $q$, yielding the next physical queue.
\end{lemma}

\begin{proof}
Each step reads one letter of $q$ and deletes two.
If $|q|$ is even, the $|q|/2$ steps consume $q$ exactly, leaving the concatenated block outputs with cut bit $p' = 0$.
If $|q|$ is odd, after $(|q|-1)/2$ steps one letter of $q$ remains; its step appends its production and deletes both that final old letter and the first newly appended letter (which exists because $w^{(1)}$ emits at least one letter).
The remaining queue is the next nominal word with its first letter removed, i.e., cut bit $p' = 1$.
In both cases $p'$ equals $p \oplus (|w| \bmod 2)$, which is the outgoing lane of the nominal sweep.
\end{proof}

\subsection{Differentiator and integrator cells}

Clocks are assembled from two 1-bit memory cells, each having an even-length state-$0$ word and an odd-length state-$1$ word:

\begin{center}
\begin{tabular}{@{}lllll@{}}
\toprule
Cell & State $0$ ($s=0$) & State $1$ ($s=1$) & Outgoing lane $y$ & Next state $s'$ \\
\midrule
Integrator     & \ttw{ezII} & \ttw{eezII} & $p \oplus s$ & $p \oplus s$ \\
Differentiator & \ttw{DDez} & \ttw{DDeze} & $p \oplus s$ & $p$ \\
\bottomrule
\end{tabular}
\end{center}

Direct inspection of positions $0,2,\ldots$ (lane $0$) and $1,3,\ldots$ (lane $1$) verifies the table:
$\ttw{ezII} \mapsto (\ttw{ezII}, \ttw{eezII})$,
$\ttw{eezII} \mapsto (\ttw{eezII}, \ttw{ezII})$,
and $\ttw{DDez}, \ttw{DDeze} \mapsto (\ttw{DDez}, \ttw{DDeze})$.
In both cells, width parity gives $y = p \oplus s$.
Thus a differentiator satisfies $y_t = p_t \oplus p_{t-1}$ and an integrator satisfies $y_t = p_t \oplus y_{t-1}$ (with zero past).
Representing a bit stream $(b_t)_{t \ge 0}$ as a formal power series $\sum_{t \ge 0} b_t X^t \in \Ftwo[[X]]$, their zero-state input--output maps are multiplication by $1+X$ and division by $1+X$, respectively.

\subsection{Clock construction and seed formula}

\begin{definition}[Seeded clock]\label{def:clock}
Let $L$ be a power of two and let $m = (m_0, \ldots, m_{L-1}) \in \Ftwo^L$ be a period-$L$ bitmask.
Define the \emph{clock word} $C(m)$ as $L$ copies of $\ttw{DDez}$ followed by $L$ integrators indexed $i = 0, \ldots, L-1$, where integrator $i$ is initialized to $\ttw{eezII}$ if $\mathrm{seed}_i = 1$ and $\ttw{ezII}$ if $\mathrm{seed}_i = 0$, with
\begin{equation}\label{eq:seed}
\mathrm{seed}_i \;=\; \bigoplus_{i \subseteq j} m_j,
\end{equation}
where $i \subseteq j$ means every set bit in the binary expansion of $i$ is also set in $j$.
\end{definition}

\begin{theorem}[Clock identity]\label{thm:clock}
For any incoming alignment stream $(p_t)_{t \ge 0}$ (including streams modified dynamically by earlier counters), the outgoing alignment of $C(m)$ at sweep $t$ is $p_t \oplus m_{t \bmod L}$.
\end{theorem}

\begin{proof}
By linearity over $\Ftwo$, decompose the output into the zero-seed response to $(p_t)$ plus the zero-input response to the initial seeds.
With all seeds zero, $L$ differentiators followed by $L$ integrators have transfer $(1+X)^L (1+X)^{-L} = 1$, passing every incoming stream $(p_t)$ unchanged.

Now set only $\mathrm{seed}_i = 1$ on zero input.
Cell $i$ emits the constant stream $1 + X + X^2 + \cdots = (1+X)^{-1}$, followed by $L - 1 - i$ integrators.
Its contribution to the output stream is
\[
(1+X)^{-(L-i)} \;=\; \frac{(1+X)^i}{(1+X)^L} \;=\; \frac{(1+X)^i}{1 + X^L},
\]
since $L$ is a power of two in characteristic $2$.
The denominator $1 + X^L$ repeats the numerator with period $L$.
Factoring $(1+X)^i = \prod_{k \in \mathrm{bits}(i)} (1 + X^{2^k})$ shows that the coefficient of $X^t$ ($0 \le t < L$) is $1$ if and only if $t \subseteq i$.
Therefore, the zero-input output at sweep $t < L$ is $\bigoplus_{t \subseteq i} \mathrm{seed}_i$.

Substituting~\eqref{eq:seed}, each $m_j$ with $t \subseteq j$ is counted once for every intermediate index $i$ satisfying $t \subseteq i \subseteq j$.
There are $2^{|j \setminus t|}$ such indices, which is odd if and only if $j = t$.
Hence the zero-input output at sweep $t$ is $m_{t \bmod L}$, and adding the zero-seed response gives $p_t \oplus m_{t \bmod L}$.
\end{proof}

\section{Counters and Guarded Instructions}

\subsection{Queue layout}

A \emph{command} consists of two sweeps: for each counter, a first-sweep \emph{sign bit} ($0$ for increment, $1$ for decrement) and a second-sweep \emph{lane bit}.
Given a nonempty even number of commands, let $L$ be the least power of two at least twice that count, and pad to $L/2$ commands using cancelling pairs of all-increment, all-decrement commands (with zero second-sweep bits) starting at an instruction boundary.

For $r$ counters, let $d_j \in \Ftwo^L$ be the resulting stream of desired lane bits at counter $j$ (interleaving both sweeps of each command), and set $d_0 = d_{r+1} = 0$.
Writing $W = \ttw{aaaeaeae}$, the starting queue for initial logical counter values $v_1, \ldots, v_r \ge 0$ is
\begin{equation}\label{eq:layout}
C(d_0 \oplus d_1)\; W^{2^{4+4v_1}}\; C(d_1 \oplus d_2)\; W^{2^{4+4v_2}}\; \cdots\; C(d_{r-1} \oplus d_r)\; W^{2^{4+4v_r}}\; C(d_r),
\end{equation}
started at cut bit $0$ (with $\oplus$ acting pointwise).
Because the clock masks telescope, Theorem~\ref{thm:clock} implies that on sweep $t$, counter $j$ receives lane $d_{j,t \bmod L}$ XOR the starting cut XOR the width parities of earlier data blocks, while the outgoing cut equals the starting cut XOR the width parities of \emph{all} data blocks.

\subsection{Exponential counters and underflow feedback}

At the start of a command, a counter with exponent $E \ge 0$ is the word $W^{2^E}$.
Reading $W = \ttw{aaaeaeae}$ on lane $0$ selects $\ttw{a,a,a,a} \to \ttw{bbbb}$; reading it on lane $1$ selects $\ttw{a,e,e,e} \to \ttw{b}$.
On the second sweep, any even run $\ttw{b}^{2k}$ regenerates $W^k$ on \emph{either} lane.
Consequently:
\begin{itemize}
\item Lane $0$ maps $W^{2^E} \to \ttw{b}^{4 \cdot 2^E} \to W^{2^{E+1}}$, incrementing $E$ by $1$.
\item Lane $1$ with $E \ge 1$ maps $W^{2^E} \to \ttw{b}^{2^E} \to W^{2^{E-1}}$, decrementing $E$ by $1$.
\item \emph{Underflow} ($E=0$, lane $1$): $W^1 \to \ttw{b}^1$. Because $|\ttw{b}| = 1$ is odd, reading $\ttw{b}$ on lane $0$ on the second sweep regenerates $W^1$ (keeping $E=0$) while toggling the cut bit for the next command.
\end{itemize}
Every first-sweep data block has even width $8 \cdot 2^E$, and on the second sweep only this singleton $\ttw{b}$ can alter downstream alignment.

\subsection{A 12-command zero-guard instruction}

At the start of a logical instruction, the cut is $0$ and each logical counter value $v \ge 0$ is encoded with exponent $E = 4 + 4v$.
To update each passive counter by a specified logical amount $a \in \{-1, 0, +1\}$ if and only if a designated tested counter equals $0$, we emit the $12$ commands in Table~\ref{tab:guard} (where $+$ denotes first-sweep lane $0$ and $-$ denotes first-sweep lane $1$).
All second-sweep lane bits are $0$, except at the tested counter in commands $6$ and $8$, where the bit is $1$.

\begin{table}[h]
\centering
\begin{tabular}{@{}lllll@{}}
\toprule
Commands & Tested & Passive ($a=+1$) & Passive ($a=0$) & Passive ($a=-1$) \\
\midrule
1--4 (remove bias)  & \ttw{----} & \ttw{+-+-} & \ttw{+-+-} & \ttw{+-+-} \\
5--6 (first probe)  & \ttw{-+}   & \ttw{+-}   & \ttw{+-}   & \ttw{-+}   \\
7--8 (second probe) & \ttw{-+}   & \ttw{+-}   & \ttw{-+}   & \ttw{-+}   \\
9--12 (restore bias)& \ttw{++++} & \ttw{+-+-} & \ttw{+-+-} & \ttw{+-+-} \\
\bottomrule
\end{tabular}
\caption{Programmed first-sweep signs for the 12-command zero-guard instruction.}
\label{tab:guard}
\end{table}

\begin{lemma}[Zero-guard invariant]\label{lem:guard}
Suppose at entry the cut is $0$, every counter has exponent $4 + 4v_j$, and whenever the tested counter has $v_{\mathrm{test}} = 0$, every passive counter with $a_j = -1$ has $v_j \ge 1$.
Then after the $12$ commands of Table~\ref{tab:guard}, the cut is again $0$, the tested counter still has exponent $4 + 4v_{\mathrm{test}}$, and each passive counter has exponent $4 + 4(v_j + a_j)$ if $v_{\mathrm{test}} = 0$ and $4 + 4v_j$ if $v_{\mathrm{test}} > 0$.
\end{lemma}

\begin{proof}
Commands 1--4 reduce the tested exponent to $4v_{\mathrm{test}}$ while leaving every passive counter unchanged (never dropping its exponent below $4$).
If $v_{\mathrm{test}} > 0$, then $4v_{\mathrm{test}} \ge 4$, so no counter underflows during probes 5--6 or 7--8; every probe executes its programmed signs, each pair cancels, and commands 9--12 restore the bias.

If $v_{\mathrm{test}} = 0$, then in each probe (5--6 and 7--8) the first $-$ hits $E=0$ and produces a singleton $\ttw{b}$.
Both the cut and the tested second-sweep bit are $0$, so the singleton is read, regenerates $W$, and flips the cut to $1$ for the second command of the probe.
Cut $1$ reverses every programmed sign in that second command: the tested counter's $+$ becomes $-$, underflowing again, and its second-sweep bit $1$ cancels cut $1$ so the second singleton is also read on lane $0$, restoring cut $0$.
Meanwhile, each passive counter executes the opposite of its second programmed sign, which equals its first programmed sign---applying that sign twice.
Summing both probes changes each passive exponent by $+4$, $0$, or $-4$ according to $a \in \{+1, 0, -1\}$, i.e., $4a$ exponent units.
Passive counters enter with exponent at least $4$ (and at least $8$ when $a=-1$ on the zero branch), preventing accidental passive underflow throughout.
\end{proof}

\section{From Counter Machines to Exact Cycles}

The reduction proceeds in two stages: a reversible front end followed by a tag compiler with a forward--backward bouncing wrapper.

\subsection{Normalized source machines}

A \emph{counter machine} has finitely many states, finitely many nonnegative integer registers, and finitely many edges, each performing one of: zero test, positivity test, increment, decrement (enabled only when positive), or no-op.
We call a counter machine \emph{standard} if it is deterministic (at most one enabled edge per configuration), has a fresh entry state (no incoming edges) and a terminal halt state (no outgoing edges), and has an enabled edge at every nonhalting configuration (reachable or not).
Standard Minsky counter machines~\cite{Minsky1967} take this form after adding a fresh entry state.

\subsection{Reversible front end}

Given a standard counter machine $M$ with $N$ edges numbered $0, \ldots, N-1$, we construct an equivalent reversible machine $\mathrm{rev}(M)$ by adding a \emph{history register} $H$ and a \emph{work register} $K$, both initialized to $0$.
Each source edge $e$ into state $q$ is replaced by a three-step gadget:
\begin{enumerate}
\item Apply the operation of $e$, move to a private state unique to $e$, and test $K = 0$.
\item Transfer $H$ into $K$ one unit at a time.
\item Enter a shared loop for $q$ at offset $\mathrm{idx}(e)$. From offset $k$, the loop adds $N + 1 - k$ to $H$ and reaches the loop head; each pass through the head moves one unit from $K$ and adds $N + 1$ to $H$, exiting to $q$ when $K = 0$.
\end{enumerate}
Net effect: $H \mapsto (N+1)H + (N + 1 - \mathrm{idx}(e))$.
The remainder modulo $N+1$ uniquely identifies the edge $e$ that entered $q$, making every step backward-deterministic.
Each gadget state has at most two incoming edges, separated by a zero test against a positivity test or increment on the same register.
The one-sided tests ($K=0$ in step~1, $H>0$ in the loop) always succeed on runs starting from $H = K = 0$.
Hence $\mathrm{rev}(M)$ has a fresh entry, terminal halt, forward- and backward-determinism, an enabled edge at every reachable nonhalting configuration, and halts if and only if $M$ halts.
(Morita~\cite[Theorem~3.1]{Morita1996} gives an alternative reversible simulation using a doubling history counter.)

\subsection{Compiling a reversible machine}

Let $R$ be a reversible machine of the form produced by $\mathrm{rev}(M)$.
Alongside the data registers, allocate a \emph{state flag} for each state ($0$ active, $1$ inactive; only entry starts at $0$), an \emph{edge latch} for each edge (all initially $1$), and a \emph{permanent zero register} for unconditional updates.
One period of the compiled tag program simulates one step of $R$ via three finite scans of 12-command zero-guard instructions:
\begin{enumerate}
\item \textbf{Select.} For each edge, temporarily add a mismatch bit to its source flag ($+1$ if operand $=0$ for a positivity test or decrement; unconditional $+1$ then $-1$ if operand $=0$ for a zero test; $0$ for increment or no-op). Guarded by the resulting source flag being $0$, decrement the edge latch, then undo the mismatch updates in reverse order with opposite signs. Exactly the unique enabled edge now has latch $0$.
\item \textbf{Apply.} For each edge, guarded by its latch being $0$, add $1$ to its source flag, subtract $1$ from its destination flag (combining updates for self-loops), and apply its data operation. Exactly one edge fires.
\item \textbf{Clear.} For each edge, run the selection test using its destination flag and the guard of its inverse operation (tests keep their predicate; inverse of increment tests positive; inverse of decrement is unconditional; no-op is self-inverse), this time \emph{incrementing} its latch. Disjoint images ensure that only the edge just taken passes, restoring all latches to $1$, all temporary flags, and the permanent zero.
\end{enumerate}
Every control register is distinct from its target, and every enabled subtraction has a positive operand by the flag, latch, and machine invariants.
Concatenating the scans, omitting empty updates, and padding to a power-of-two period yields a queue whose period boundaries track $R$'s trajectory by induction.

\subsection{Halting as exact queue repetition}

\begin{proof}[Proof of Theorem~\ref{thm:main}]
Apply the reversible front end to obtain $R = \mathrm{rev}(M)$, then construct a bouncing reversible machine $B(R)$ with forward and backward copies of every state: keep original edges in the forward copy, reverse each edge with its inverse operation in the backward copy, and add two no-op turnarounds (forward halt $\to$ backward halt, backward entry $\to$ forward entry).
Missing incoming edges at entry and outgoing edges at halt keep $B(R)$ reversible.
Compile $B(R)$ starting at forward entry.

If $M$ halts, then $R$ halts, so $B(R)$ bounces indefinitely along a finite path between entry and halt.
Register values at period boundaries are bounded.
Within a period, each exponent can grow by at most the fixed number of commands, and each clock cell has width $4$ or $5$; intermediate queues during a sweep consist of an old suffix plus a prefix of the bounded next word, so physical queue lengths are globally bounded.
By Lemma~\ref{lem:sweep}, the leading clock sustains every sweep.
Finitely many words of bounded length over $\Sigma$ force exact repetition.

If $M$ does not halt, then $R$ runs forever in the forward copy.
Its trajectory never repeats a configuration: by backward-determinism, any repetition would pull back to a predecessor of the fresh entry state, which does not exist.
Finitely many control states therefore force unbounded register values at period boundaries, hence unbounded physical queue lengths.
Exact repetition at any tag step would trap the deterministic tag system in a periodic tail of bounded length, a contradiction.
\end{proof}

\section{Formal Verification in Lean 4}

The entire reduction is formalized in two Lean~4 files (checked with Lean~4.19 and its standard library, using only the standard axioms \ttw{propext}, \ttw{Classical.choice}, and \ttw{Quot.sound}, with no \ttw{sorryAx}):
\begin{itemize}
\item \ttw{SixSymbol.lean} formalizes the 2-tag system~\eqref{eq:productions}, clock identities, zero-guard instructions, reversible compiler, bouncing wrapper, and \ttw{SixSymbol.halt\_iff\_exact\_cycle}.
\item \ttw{ReversibleCounter.lean} imports \ttw{SixSymbol.lean} unchanged, formalizes the reversible front end \ttw{rev}, discharges its reversibility and liveness obligations, and proves the end-to-end theorem (\ttw{ReversibleCounter.\allowbreak halt\_iff\_exact\_cycle}):
\begin{align*}
(\exists n,\; (m.\mathrm{path}\;\mathrm{input}\;n).\mathrm{pc} = m.\mathrm{halt})
&\;\longleftrightarrow\; \\
&\mathrm{ExactCycle}\bigl(\mathrm{encode}(\mathrm{rev}\;m\;\mathrm{std.fresh})(\mathrm{liftInput}\;\mathrm{input})\bigr).
\end{align*}
\end{itemize}

\begin{table}[h]
\centering
\begin{tabular}{@{}lp{0.58\linewidth}@{}}
\toprule
Theorem & Formalized statement \\
\midrule
\ttw{seeded\_clock}, \ttw{layout\_sweep} & Seeded clocks emit the target periodic mask on either lane; whole-layout sweeps match literal tag steps, including odd-boundary crossings. \\
\ttw{gate\_realized}, \ttw{compiled\_periods} & Two-sweep commands realize the 12-command zero guard on a register bank; padded programs execute successive periods on one queue. \\
\ttw{compiled\_repeats\_iff} & Exact repetition holds iff period-boundary register files range over a finite set (even if clock states vary between visits). \\
\ttw{one\_period}, \ttw{checkpoint\_edge} & Select, Apply, and Clear realize the enabled forward/inverse step; turnarounds realize the bouncing trajectory. \\
\ttw{halt\_iff\_exact\_cycle} & Halting $\iff$ exact queue repetition for reversible live machines. \\
\ttw{rev\_reversible}, \ttw{rev\_live} & $\mathrm{rev}(M)$ is reversible and live on $H=K=0$ inputs. \\
\ttw{rev\_of\_halt}, \ttw{halt\_of\_rev} & $\mathrm{rev}(M)$ halts iff $M$ halts, yielding the end-to-end equivalence. \\
\bottomrule
\end{tabular}
\caption{Main verified theorems in \ttw{SixSymbol.lean} (top five rows) and \ttw{ReversibleCounter.lean} (bottom three rows).}
\label{tab:lean}
\end{table}

The sole unformalized premise is Minsky's theorem~\cite{Minsky1967} together with the routine normalization of a counter machine to a fresh entry and terminal halt.
An executable Python reference implementation of the reversible compiler is provided in \ttw{parity\_program.py}.

\section{Examples}

\paragraph{Doubling a tile.}
Starting from $W = \ttw{aaaeaeae}$, six literal 2-tag steps double the block:
\[
\ttw{aaaeaeae} \to \ttw{aeaeaeb} \to \ttw{aeaebb} \to \ttw{aebbb} \to \ttw{bbbb}
\to \ttw{bbaaaeaeae} \to \ttw{aaaeaeaeaaaeaeae}.
\]

\paragraph{FizzBuzz.}
A single 2,128-letter starting queue generates FizzBuzz indefinitely (see \ttw{six\_symbol\_\allowbreak fizzbuzz.py}).
Each period has $16$ slots of $4$ sweeps ($L = 64$).
Numbering sweeps $t = 4(s-1) + u$ ($1 \le s \le 16$, $0 \le u < 4$), define $F_t = 1$ at offset $0$ of slots $3,6,9,12,15$; $B_t = 1$ at offset $0$ of slots $5,10,15$; and $A_t = 1$ only at slot $16$, offset $2$.
The initial queue is
\[
C(F)\; C(F \oplus B)\; C(B \oplus A)\; W\; C(A),
\]
whose four 64-bit integrator seed vectors (cell $0$ as LSB) in hexadecimal are
\begin{center}
\ttw{0101101101101101},\quad \ttw{0000110100101010},\quad \ttw{5454544554555444},\quad \ttw{5555555555555555},
\end{center}
giving lengths $522 + 518 + 536 + 8 + 544 = 2{,}128$.
In slots $1$--$15$, the first two clocks signal Fizz and Buzz while the data block doubles twice per slot (holding $4^n$ copies of $W$ after output $n$); slot $16$ doubles once and halves once, keeping the data run even on all sweeps so clock~4 restores cut $0$ by induction.

\paragraph{Bouncing machine.}
Compiling a single-edge machine $(\ttw{s} \xrightarrow{+X} \ttw{h})$ produces a 512-sweep period.
The source halts after period~$1$, while the tag queue bounces $\ttw{s}(0) \to \ttw{h}(1) \to \ttw{h}(1) \to \ttw{s}(0) \to \ttw{s}(0) \to \ttw{h}(1)$, matching identically (including every clock cell) between checkpoints~$1$ and~$5$, $2{,}048$ sweeps apart.

\section{Minimality and Five-Symbol Perspectives}

Fix deletion number $2$, a single production table, finite starting queues, and exact queue repetition as the observable.
Since binary ($2$-symbol) $2$-tag systems are decidable~\cite{DeMol2010} and Theorem~\ref{thm:main} establishes universality at $6$ symbols, the remaining open frontier for universal $2$-tag systems is alphabet sizes $3$, $4$, and $5$ (with $3$ symbols already hosting De~Mol's Collatz system $\ttw{a} \to \ttw{bc}$, $\ttw{b} \to \ttw{a}$, $\ttw{c} \to \ttw{aaa}$~\cite{DeMol2008}).

\begin{proposition}[Three-letter clock obstruction]\label{prop:obstruction}
Reserve two symbols for data and require a disjoint, closed clock alphabet (containing an eraser $\ttw{e}$) in which each clock cell uses two fixed state words of opposite length parity that regenerate the prescribed state word in one sweep on either alignment.
Then no three-letter clock alphabet can implement both an integrator and a differentiator under a single table; hence at least $4 + 2 = 6$ symbols are required for this clock architecture.
\end{proposition}

\begin{proof}[Proof sketch]
Let $A,B$ be the two non-eraser clock letters, $M$ the $2 \times 2$ matrix counting $(A,B)$ in their appendants, and $c_0, c_1$ the count vectors of a cell's even and odd state words.
Summing both lanes gives $M c_0 = M c_1 = c_0 + c_1$.
If $M$ were invertible, $c_0 = c_1$ would force equal output lengths on both lanes, contradicting opposite state parities.
Thus $M$ is singular with eigenvalue $2$, leaving two cases:
(i) all entries of $M$ are $1$, or
(ii) the diagonal entries are $2$ and $0$.
In both cases, every state word parses into $n > 0$ copies of the generating appendant(s), and opposite state parities force each copy to place its live letters in opposite lanes.
Every copy then obeys the same linear difference transfer in $\{1,\, 1+X,\, (1+X)^{-1}\}$, so positive powers of a single transfer cannot realize both $1+X$ and $(1+X)^{-1}$.
Full details are given in \ttw{clock\_obstruction.tex}.
\end{proof}

Two five-symbol variations highlight what remains open beyond Proposition~\ref{prop:obstruction}:
\begin{enumerate}
\item \emph{Removing $\ttw{D}$ (integrator-only horizon compiler).} Dropping $\ttw{D}$ yields a five-symbol system that matches the six-symbol counters up to any chosen finite horizon $K$ (a power of two divisible by the period) by repeating each mask to length $K$ and seeding $K$ integrators, whose transfer $(1+X)^{-K} = (1+X^K)^{-1}$ copies inputs faithfully through sweep $K-1$. For the zero-guard example from $(0,0)$ with $K=32$, counter $Y$ agrees through sweep $32$ and diverges at sweep $34$ ($512$ vs.\ $128$ copies of $W$).
\item \emph{Data-reloaded control.} A separate five-symbol table (\ttw{five-symbol-tag.md}) uses data itself to reload control, already demonstrating a second load from a retained payload; whether it supports recurring simulation from live output remains open.
\end{enumerate}
Settling whether six symbols are minimal in general requires ruling out data-reloaded control and all other five-symbol tables and encodings.

\begin{thebibliography}{9}

\bibitem{ais523UT19}
ais523.
\newblock UT19 (a 19-symbol universal 2-tag system).
\newblock Esolang Wiki, 2023.
\newblock \url{https://esolangs.org/wiki/UT19}.

\bibitem{CockeMinsky1964}
J.~Cocke and M.~Minsky.
\newblock Universality of tag systems with $P=2$.
\newblock \emph{Journal of the ACM}, 11(1):15--20, 1964.

\bibitem{Cook2004}
M.~Cook.
\newblock Universality in elementary cellular automata.
\newblock \emph{Complex Systems}, 15(1):1--40, 2004.

\bibitem{DeMol2008}
L.~De~Mol.
\newblock Tag systems and Collatz-like functions.
\newblock \emph{Theoretical Computer Science}, 390(1):92--101, 2008.

\bibitem{DeMol2010}
L.~De~Mol.
\newblock Solvability of the halting and reachability problem for binary 2-tag systems.
\newblock \emph{Fundamenta Informaticae}, 2010.
\newblock (See also arXiv:0906.3329.)

\bibitem{Minsky1961}
M.~L.~Minsky.
\newblock Recursive unsolvability of Post's problem of ``Tag'' and other topics in theory of Turing machines.
\newblock \emph{Annals of Mathematics}, 74(3):437--455, 1961.

\bibitem{Minsky1967}
M.~L.~Minsky.
\newblock \emph{Computation: Finite and Infinite Machines}.
\newblock Prentice-Hall, 1967.

\bibitem{Morita1996}
K.~Morita.
\newblock Universality of a reversible two-counter machine.
\newblock \emph{Theoretical Computer Science}, 168(2):303--320, 1996.

\bibitem{NearyWoods2006}
T.~Neary and D.~Woods.
\newblock P-completeness of cellular automaton Rule 110 (and polynomial-time simulation of 2-tag systems).
\newblock In \emph{ICALP 2006}, LNCS 4051, pp.~132--143, Springer, 2006.

\bibitem{NearyWoods2009}
T.~Neary and D.~Woods.
\newblock Four small universal Turing machines.
\newblock \emph{Fundamenta Informaticae}, 91(1):123--144, 2009.

\bibitem{Post1943}
E.~L.~Post.
\newblock Formal reductions of the general combinatorial decision problem.
\newblock \emph{American Journal of Mathematics}, 65(2):197--215, 1943.

\bibitem{Rogozhin1996}
Y.~Rogozhin.
\newblock Small universal Turing machines.
\newblock \emph{Theoretical Computer Science}, 168(2):215--240, 1996.

\bibitem{Wang1963}
H.~Wang.
\newblock Tag systems and lag systems.
\newblock \emph{Mathematische Annalen}, 152(1):65--74, 1963.

\end{thebibliography}

\end{document}
